\documentclass[%
	english,%
	 a4paper,%
	%
	%
	%
	fleqn]{article}
% This is ugly, but if I have it spaced out, pandoc doesn't like it much.

% Set your local directory---*do not* include a slash after the directory
\def\localDir{/home/pines/.pandoc/quetzalcoatl}

% DO NOT include a slash after your fontDir is set
\def\fontDir{\localDir/fonts}
\def\packageDir{\localDir/packages}

% I used to use the `silence` package here, but I found that you can get rid
% of the annoying popups by simply \emph{not} including the `\ProvidesPackage`
% command.

%\usepackage{authblk}

\usepackage{amsmath}
\usepackage{amssymb}
\usepackage{amsthm}
\usepackage[makeroom]{cancel}
\usepackage{xparse}
\usepackage{xstring}
\usepackage{xtemplate}
\usepackage{xspace}
\usepackage{ragged2e}
\usepackage[palette=default]{\packageDir/mycolors}
\usepackage{\packageDir/subtitle}


% Physics stuff
% \usepackage{braket}
\usepackage[spaced=0]{diffcoeff}
\usepackage{linop}

% Change the diffcoeff "d" to be upright.
    \difdef{ f, s, c, l }{}
        {
            op-symbol = d ,
            op-order-nudge = 1 mu
        }


\usepackage{\packageDir/mymacros}

\usepackage{enumerate}
% Keep pandoc from getting mad
\providecommand{\tightlist}{%
  \setlength{\itemsep}{0pt}\setlength{\parskip}{0pt}}

\renewcommand{\labelenumii}{\arabic{enumi}.\arabic{enumii}.}
\renewcommand{\labelenumiii}{\arabic{enumi}.\arabic{enumii}.\arabic{enumiii}.}
\renewcommand{\labelenumiv}{%
    \arabic{enumi}.\arabic{enumii}.\arabic{enumiii}.\arabic{enumiv}.}

\usepackage{calc}


\usepackage{\packageDir/mathParenSpacing}

\usepackage{soul}

\usepackage{siunitx}

\usepackage{fontspec}
\usepackage{unicode-math}
\usepackage[noindentafter]{titlesec}

\usepackage[auto=true]{microtype}
\SetTracking{encoding={*}, shape=sc}{500}

    \usepackage[style=regular]{\packageDir/stylesets}



% %     	\maketitle
%         % 

	\title{Abscence of diffusion in certain random lattices}

    \subtitle{testing out titles and stuff}

	\author{P.W. Anderson}


	\date{1957-10-10}

    \newcommand{\QuetzalAbstract}{{This paper presents a simple model
for such processes as spin diffusion or conduction in the ``impurity
band.'' These processes involve transport in a lattice which is in some
sense random, and in them diffusion is expected to take place via
quantum jumps between localized sites. In this simple model the
essential randomness is introduced by requiring the energy to vary
randomly from site to site. It is shown that at low enough densities no
diffusion at all can take place, and the criteria for transport to occur
are given.}}

\usepackage[runin]{abstract}
\setlength{\absparindent}{0 pt}
\renewcommand{\abstractnamefont}{\sffamily\small\bfseries}
    \renewcommand\abstractname{\sffamily Abstract}

% Title formatting
    % \usepackage[runin]{abstract}

\usepackage{\packageDir/title}




% Set itemize symbols
\renewcommand{\labelitemi}{\symbolFont{▪}}
\renewcommand{\labelitemii}{\symbolFont{▹}}
\renewcommand{\labelitemiii}{\symbolFont{•}}
\renewcommand{\labelitemiv}{\symbolFont{⟶}}

% \usepackage{indentfirst}




	\usepackage{polyglossia}
	\setmainlanguage{english}




\usepackage{setspace}
	\setstretch{1.25}


\usepackage{longtable,booktabs,array}
\usepackage{multirow}

\usepackage{afterpage}


\setlength\heavyrulewidth{0.3ex}      % thickness of \toprule, \bottomrule
\renewcommand{\arraystretch}{1.3}     % spacing (padding)


% \usepackage{bm}

\usepackage{upquote}

\usepackage[hmargin=2in, marginparsep=0.0 in]{geometry}


\usepackage{fancyhdr}
        \pagestyle{plain}



% Taken from Eisvogel, ln. 312--339
% https://github.com/Wandmalfarbe/pandoc-latex-template/blob/master/eisvogel.tex





% Also taken from Eisvogel

\usepackage{graphicx}
\DeclareGraphicsExtensions{.png,.pdf}

% Useful---I like it.
\usepackage{subcaption}


\makeatletter
\def\maxwidth{\ifdim\Gin@nat@width>\linewidth\linewidth\else\Gin@nat@width\fi}
\def\maxheight{\ifdim\Gin@nat@height>\textheight\textheight\else\Gin@nat@height\fi}
\makeatother
% Scale images if necessary, so that they will not overflow the page
% margins by default, and it is still possible to overwrite the defaults
% using explicit options in \includegraphics[width, height, ...]{}
\setkeys{Gin}{width=\maxwidth,height=\maxheight,keepaspectratio}
%% Set default figure placement to htbp
% Set default figure placement to t
\makeatletter
%\def\fps@figure{htbp}
\def\fps@figure{t}
\makeatother

% Change figure title
\usepackage[figurename={Figure}, font={small}, labelfont={sf, bf},%
labelsep={period}, justification={raggedright}]{caption}
% \DeclareCaptionStyle{marginfigure}{labelfont={color=mymaroon}}



\usepackage[unicode=true]{hyperref}


\hypersetup{breaklinks=true,
            bookmarks=true,
            pdfauthor={P.W. Anderson},
            pdftitle={Abscence of diffusion in certain random lattices},
            colorlinks=true,
            citecolor= black ,
	        urlcolor=urlColor,
            linkcolor= black ,
            pdfborder={0 0 0},
}
%\def\UrlBreaks{\do\/\do-}
\urlstyle{same}  % don't use monospace font for urls


%\definecolor{blockquote-border}{RGB}{221,221,221}
% redefine footnotes to include character

% It's time to get funky!
\let\periodGroup=.
%\let\commaGroup=,
%\newcommand{\superUrlMark}{
    %\@ifnextchar.{\rlap{°}}{°}
%}

% make something that automatically kerns a urlMark. I didn't want to switch to
% luaTeX, since XeTeX is slow enough as is.
%
% Check a `\url` command, and see if the character next to the hyperlink is a
% period, or see if the last character of a hyperlink is a period or comma.
%
%\dots\stackengine{0pt}{.}{}{O}{l}{F}{F}{L}


	\let\orighref\href
	%\renewcommand{\href}[2]{\orighref{#1}{#2\color{urlMarkColor}%
	%\rlap{\symbolFont{°}}}} % I wish this would work properly...
	%\renewcommand{\href}[2]{\orighref{#1}{#2\color{urlMarkColor}%
	%\symbolFont{°}}}
	%\renewcommand{\href}[2]{\orighref{#1}{#2\color{urlMarkColor}%
        %°}} % MAYBE COME BACK TO THIS
	%\renewcommand{\href}[2]{\orighref{#1}{#2\color{urlMarkColor}%
    %\superUrlMark}}



\setlength{\emergencystretch}{3em}  % prevent overfull lines

% No clue why I put in numberSections twice
%%\setcounter{secnumdepth}{0}
%



\usepackage{mdframed}
\newmdenv[rightline=false, leftline=false, linecolor=ruleColor]{abstractBox}

\newmdenv[rightline=false,bottomline=false,topline=false,linewidth=3pt,linecolor=blockquote-border,skipabove=1em,font=\sffamily\small]{customblockquote}
\renewenvironment{quote}{\begin{customblockquote}\vspace{-0.5cm}{\rightmargin=0em\leftmargin=0em}%
\item\relax\color{blockquote-text}\ignorespaces}{\unskip\unskip\end{customblockquote}}

	\setlength{\parindent}{1.25em}



		\setlength\RaggedRightParindent{1.25em}
	

\newcommand{\dinkus}{\begin{center}* * *\end{center}}

\widowpenalty=2000
\clubpenalty=2000

	\setlength{\parskip}{0cm plus 1mm minus 1mm}

\usepackage{lipsum}

\usepackage{datetime2}

\usepackage{tikz}


\usetikzlibrary{positioning}
\usetikzlibrary{shapes.geometric}

\usepackage{multicol}

% Taken from the default template
\usepackage[]{biblatex}
\addbibresource{ref.bib}


% foo
%\if\texttt1
	%\let\emdash{\emdash\emdash}
%\fi

% !!!ALWAYS MAKE SURE THIS IS LOADED *LAST*!!!


\begin{document}


    	\maketitle
        



%   % Maybe I could get rid of this?
%     \let\QuetzalAbstract\relax
% 


% 

    %    %\renewcommand\abstractname{\sffamily Abstract.}
    %
%\setlength{\absparindent}{0.51em}

% possible spaghetti incoming
% % \setlength{\absparindent}{0.51em}


%	\begin{abstract}
%	\begin{abstractBox} % sometimes we have happy little accidents
%		abstract  % surround this in dollar signs when you want it to work with pandoc again
%	\end{abstractBox}
%	\end{abstract}


% \saythanks

 





% TODO: work on developing this more



%%%\SetBgContents{\rule{.4pt}{\textheight + 2\headsep}}
%\pagestyle{fancy}
\section{Introduction}\label{introduction}

A number of physical phenomena seem to involve quantum-mechanical
motion, without any particular thermal activation, among sites at which
the mobile entities (spins or electrons, for example) may be localized.
The clearest case is that of spin diffusion''
\autocite{bloembergen_1949,portis_1956}; another might be the so-called
impurity band conduction at low concentrations of impurities. In such
situations we suspect that transport occurs not by motion of free
carriers (or spin waves), scattered as they move through a medium, but
in some sense by quantum-mechanical jumps of the mobile entities from
site to site. A second common feature of these phenomena is randomness:
random spacings of impurities, random interactions with the
``atmosphere'' of other impurities, random arrangements of electronic or
nuclear spins, etc.

Our eventual purpose in this work will be to lay the foundation for a
quantum-mechanical theory of transport problems of this type. Therefore,
we must start with simple theoretical models rather than with the
complicated experimental situations on spin diffusion or impurity
conduction. In this paper, in fact, we attempt only to construct, for
such a system, the simplest model we can think of which still has some
expectation of representing a real physical situation reasonably well,
and to prove a theorem about the model. The theorem is that at
sufficiently low densities, transport does not take place; the exact
wave functions are localized in a small region of space. We also obtain
a fairly good estimate of the critical density at which the theorem
fails. An additional criterion is that the forces be of sufficiently
short range---actually, falling off as \(r \rightarrow \infty\) faster
than \(1/r^3\)---and we derive a rough estimate of the rate of transport
in the \(V \propto 1/r^3\) case.

\ldots{}

\section{Summary of the reasoning}\label{summary-of-the-reasoning}

Since the mathematical development is fairly complicated and involves
lengthy consideration of each of a number of points, we should like to
summarize the reasoning rather fully in this section, leaving the proofs
and details to later sections. First, then, let us set up the simple
model which we study. The equation for the time-dependence of the
probability amplitude \(a\), that a particle is on the site \(j\) is:
\begin{equation}\phantomsection\label{eq:1}{i \, \dot{a}_j = E_j \, a_j + \sum_{k \neq j} V_{j k} \, a_k \,.}\end{equation}
Here we measure energies in frequency units, so we can set
\(\hbar = 1\). Equation \ref{eq:1} simply restates the assumptions about
the model made in the Introduction. We study the Laplace transform of
the equation (\ref{eq:1}): let
\begin{equation}\phantomsection\label{eq:2}{ f_j(s) = \int_{0}^{\infty} e^{-s \, t} a_j(t) \dl{t} \,, }\end{equation}
and then
\begin{equation}\phantomsection\label{eq:3}{ i \left[ s \, f_j(s) - a_j(0) \right] = E_j \, f_j + \sum_{k \neq j} V_{j k} \, f_k \,.}\end{equation}
The variables \(s\) must be as an arbitrary complex frequency with
positive or zero real part.

\texttt{Now\ this\ is\ being\ written\ on\ a\ typewriter.}






\printbibliography

\end{document}
